\section{Annexe}
\subsection{Résolution intégrale des facteurs de vue}
\label{method_integrale}
On considère le facteur de vue entre deux faces opposées d'un cube unité : la face $z=0$ (surface $A_1$) et la face $z=1$ (surface $A_2$). Le facteur de vue est défini de manière générale par :
\begin{equation}
F_{16}=\frac{1}{A_1}\int_{A_1}\int_{A_2}\frac{\cos\theta_1 \cos\theta_2}{\pi R^2}\, dA_2\, dA_1.
\end{equation}

Dans le cas du cube unité, on a $A_1=A_2=1$. Les deux faces étant parallèles, on paramètre les surfaces par :
\begin{align}
P_1=(x,y,0), \quad (x,y)\in[0,1]^2,\\
P_2=(x',y',1), \quad (x',y')\in[0,1]^2.
\end{align}

La distance entre deux éléments différentiels s'écrit :
\begin{equation}
R^2=(x'-x)^2+(y'-y)^2+1.
\end{equation}

\paragraph{Calcul des cosinus directionnels}

Les normales aux surfaces sont :
\begin{equation}
\mathbf n_1=(0,0,1), \qquad \mathbf n_2=(0,0,-1).
\end{equation}

Le vecteur reliant les deux points est :
\begin{equation}
\mathbf r=(x'-x,\;y'-y,\;1), \qquad R=\|\mathbf r\|.
\end{equation}

On en déduit :
\begin{equation}
\cos\theta_1=\frac{1}{R}, \qquad \cos\theta_2=\frac{1}{R}, \qquad \Rightarrow \qquad \cos\theta_1\cos\theta_2=\frac{1}{R^2}.
\end{equation}

Ainsi, le noyau géométrique devient :
\begin{equation}
\frac{\cos\theta_1 \cos\theta_2}{\pi R^2}=\frac{1}{\pi R^4}=\frac{1}{\pi\left[1+(x'-x)^2+(y'-y)^2\right]^2}.
\end{equation}

\paragraph{Réduction par changement de variables et convolution}

L'intégrande ne dépend que des différences de coordonnées. On introduit :
\begin{equation}
u=x'-x, \qquad v=y'-y.
\end{equation}

On utilise la propriété de convolution :
\begin{equation}
\int_0^1\int_0^1 f(x'-x)\,dx\,dx'=\int_{-1}^{1}(1-|u|)\,f(u)\,du.
\end{equation}

En deux dimensions, cela donne :
\begin{equation}
F_{16}=\frac{1}{\pi}\int_{-1}^{1}\int_{-1}^{1}\frac{(1-|u|)(1-|v|)}{(1+u^2+v^2)^2}\,du\,dv.
\end{equation}

Par symétrie :
\begin{equation}
F_{16}=\frac{4}{\pi}\int_0^1\int_0^1\frac{(1-u)(1-v)}{(1+u^2+v^2)^2}\,du\,dv.
\end{equation}

\paragraph{Interprétation}

Le facteur $(1-|u|)(1-|v|)$ représente la densité de paires de points des deux surfaces ayant une séparation relative $(u,v)$. Cette transformation permet de passer d'une intégration sur toutes les paires de points à une intégration sur les distances relatives, ce qui est une étape clé en transfert radiatif et en géométrie intégrale.
\begin{equation}
F_{16}=\frac{4}{\pi}\int_0^1\int_0^1\frac{(1-u)(1-v)}{(1+u^2+v^2)^2}\,du\,dv.
\end{equation}

\paragraph{Introduction d'une fonction auxiliaire}

Afin de traiter cette intégrale, on introduit la fonction paramétrée

\begin{equation}
J(a)=\int_0^1\int_0^1\frac{(1-u)(1-v)}{a+u^2+v^2},du,dv,\qquad a>0.
\end{equation}

Cette définition conserve le facteur de pondération géométrique $(1-u)(1-v)$ issu de la réduction par convolution. En dérivant sous le signe intégral, on obtient

\begin{equation}
\frac{dJ}{da}=-\int_0^1\int_0^1\frac{(1-u)(1-v)}{(a+u^2+v^2)^2},du,dv.
\end{equation}

Par conséquent,

\begin{equation}
\int_0^1\int_0^1\frac{(1-u)(1-v)}{(1+u^2+v^2)^2},du,dv=-\left.\frac{dJ}{da}\right|_{a=1}.
\end{equation}

Le facteur de vue s'écrit alors

\begin{equation}
F_{16}=-\frac{4}{\pi}\left.\frac{dJ}{da}\right|_{a=1}.
\end{equation}

\paragraph{Décomposition de l'intégrale}

On développe le facteur de pondération :

\begin{equation}
(1-u)(1-v)=1-u-v+uv.
\end{equation}

Ainsi,

\begin{equation}
F_{16}=\frac{4}{\pi}(A-2B+C),
\end{equation}

où

\begin{equation}
A=\int_0^1\int_0^1\frac{du,dv}{(1+u^2+v^2)^2},
\end{equation}

\begin{equation}
B=\int_0^1\int_0^1\frac{u,du,dv}{(1+u^2+v^2)^2},
\end{equation}

et

\begin{equation}
C=\int_0^1\int_0^1\frac{uv,du,dv}{(1+u^2+v^2)^2}.
\end{equation}

Le calcul de ces intégrales conduit finalement à l'expression analytique classique du facteur de vue entre deux faces opposées d'un cube unité :

\begin{equation}
F_{16}=\frac{2}{\pi} \left[ \ln(\frac{2}{\sqrt{3}}) +2\sqrt{2}arctan(\frac{1}{\sqrt{2}})-\frac{\pi}{2} \right]
\end{equation}

Numériquement,

\begin{equation}
F_{16}\simeq 0.1998249.
\end{equation}